% 347_Gell_Mann_Nishijima_Galois_De_ch.tex
% Johann Pascher, 4. September 2026

\providecommand{\BM}{\textbf{[B]}}
\providecommand{\KM}{\textbf{[K]}}
\providecommand{\SM}{\textbf{[S]}}

\phantomsection
\addcontentsline{toc}{chapter}{Gell-Mann-Nishijima aus der Galois-Struktur}

\section*{Statuszeichen}
\begin{center}
\begin{tabular}{@{}lp{10.5cm}@{}}
\textbf{[K]} & \emph{Konfirmiert} --- algebraisch vollständig. \\[4pt]
\textbf{[B]} & \emph{Belegt} --- durch Prüfskript numerisch gesichert. \\[4pt]
\textbf{[S]} & \emph{Spekulativ} --- Vermutung, noch offen. \\
\end{tabular}
\end{center}

\section*{Abstract}

Untersucht wird, wie weit die Gell-Mann-Nishijima-Relation $Q = I_3 + Y/2$ aus
zwei bereits etablierten Galois-Eigenschaften folgt. (1)~Der Isospin $I_3$ folgt
aus der Su/Sd-Trennung (Dok.~344): $I_3 = \pm\tfrac{1}{2}$ \BM.
(2)~Die Hyperladung $Y$ ist aus dem QR/NQR-Bit mod~13 (Dok.~346) allein
\emph{nicht} festgelegt: geladene Leptonen und Quarks sind beide NQR, und das
Leptondublett $(\nu,e)_L$ liegt auf beiden Seiten des Bits (Satz~A, \SM).
Mit den SM-Werten $Y=-1$ (Leptondublett) und $Y=+\tfrac{1}{3}$ (Quarkdublett)
liefert die Relation die vier Fermionladungen $\{0, -1, +\tfrac{2}{3}, -\tfrac{1}{3}\}$;
Satz~E aus Dok.~346 ist damit für $I_3$ geschlossen, für $Y$ offen \SM.

\section{Ausgangslage}

Dok.~346 hat die Fermionklassifikation aus zwei Galois-Eigenschaften etabliert:
QR/NQR-Bit (elektrische Neutralität) und $N\bmod3$ (Farbladung). Als Satz~E~[\SM]
blieb die Frage: Folgt auch $Q = I_3 + Y/2$ aus der Galois-Struktur?

\section{Satz A --- Hyperladung aus QR/NQR}

\begin{theorem}[Hyperladung $Y$ ist durch das QR/NQR-Bit allein nicht festgelegt \SM]
Nach Dok.~346 (Satz~B und Orbit-Tabelle) gilt: QR-Orbits $O_1=\{1,3,9\}$ $=$ Neutrinos,
$O_3=\{4,10,12\}$ $=$ Antineutrinos; NQR-Orbits $O_4=\{7,8,11\}$ $=$ geladene Leptonen,
$O_2=\{2,5,6\}$ $=$ Quarks. Es gibt keine Funktion $Y = f(\mathrm{NQR}_{13})$, die die
SM-Werte $Y=-1$ (Leptondublett $(\nu,e)_L$) und $Y=+\tfrac{1}{3}$ (Quarkdublett $(u,d)_L$)
reproduziert. Die Zuordnung von $Y$ muss zusätzlich $N\bmod3$ und die
Teilchen/Antiteilchen-Unterscheidung einbeziehen; sie ist offen.
\end{theorem}

\begin{proof}
Angenommen $Y = f(\mathrm{NQR}_{13})$.
Neutrino ($O_1$, QR): $f(0) = -1$.
Elektron ($O_4$, NQR), im selben Dublett wie das Neutrino: $f(1) = -1$.
up-Quark ($O_2$, NQR): $f(1) = +\tfrac{1}{3}$ --- Widerspruch zu $f(1)=-1$.
Geladene Leptonen und Quarks liegen beide auf der NQR-Seite, das Leptondublett
$(\nu,e)_L$ auf beiden Seiten des Bits; das Bit trennt also weder Leptonen von
Quarks noch die Dublettpartner zusammen. Zusätzlich trägt der QR-Orbit $O_3$
(Antineutrinos) das umgekehrte Vorzeichen von $Y$ wie $O_1$, obwohl beide QR sind. \qed
\end{proof}

\textbf{Verbindung zu Dok.~346.} Das NQR-Bit ist exakt das Legendre-Symbol
$\left(\tfrac{k}{13}\right) = -1$ --- dieselbe Eigenschaft, die in Dok.~346 die
elektrische Neutralität klassifiziert \BM. Die Hyperladung folgt daraus nicht (Satz~A) \SM.

\section{Satz B --- Isospin aus Su/Sd-Trennung}

\begin{theorem}[Isospin $I_3$ aus Su/Sd-Trennung \BM]
\[
  I_3 = \begin{cases}
    +\tfrac{1}{2} & \text{Su-Sektor (Dok.~344: }F \cdot jE7 = +F\text{)}\\
    -\tfrac{1}{2} & \text{Sd-Sektor (Dok.~344: }F \cdot jE7 = -F\text{)}
  \end{cases}
\]
\end{theorem}

Die Su/Sd-Trennung durch den jE7-Projektor (Dok.~344, Satz~E) ist exakt die
Dublett-Aufspaltung des schwachen Isospins. Das ist keine neue Annahme ---
der jE7-Projektor wurde bereits in der GALG/Furey-Konstruktion etabliert.

\section{Satz C --- Gell-Mann-Nishijima mit Isospin aus Galois}

\begin{theorem}[Gell-Mann-Nishijima mit $I_3$ aus Satz~B und SM-Hyperladungen \SM]
\[
  Q = I_3 + \frac{Y}{2}
\]
mit $I_3$ aus Satz~B und den SM-Werten $Y$ (nicht aus Satz~A ableitbar) ergibt alle Fermionladungen exakt:

\begin{center}
\begin{tabular}{lllll}
\toprule
Fermion & $I_3$ & $Y$ & $Q = I_3 + Y/2$ & Exp. \\
\midrule
Neutrino & $+\tfrac{1}{2}$ & $-1$ & $0$ & $0$ \checkmark \\
\multicolumn{5}{l}{\small Vergleichswerte: PDG 2022 \cite{PDG2022}; alle vier Ladungen exakt reproduziert.}\\
Elektron & $-\tfrac{1}{2}$ & $-1$ & $-1$ & $-1$ \checkmark \\
up-Quark & $+\tfrac{1}{2}$ & $+\tfrac{1}{3}$ & $+\tfrac{2}{3}$ & $+\tfrac{2}{3}$ \checkmark \\
dn-Quark & $-\tfrac{1}{2}$ & $+\tfrac{1}{3}$ & $-\tfrac{1}{3}$ & $-\tfrac{1}{3}$ \checkmark \\
\bottomrule
\end{tabular}
\end{center}
\end{theorem}

\textbf{Alle vier Werte stimmen \BM; galoistheoretisch hergeleitet ist dabei nur $I_3$, die Hyperladung $Y$ ist als SM-Wert eingesetzt \SM.}

\section{Satz D --- Hyperladung und Legendre-Symbol (Orbit-Ebene)}

Legendre-Symbol, Teilchenzuordnung nach Dok.~346 und SM-Hyperladung (linkshändiges
Dublett bzw.\ dessen Antiteilchen) auf Orbit-Ebene:

\begin{center}
\begin{tabular}{llll}
\toprule
Orbit & $\left(\tfrac{k}{13}\right)$ & Teilchen (Dok.~346) & $Y$ (SM) \\
\midrule
$O_1=\{1,3,9\}$ & $+1$ (QR) & $\nu$ & $-1$ \\
$O_3=\{4,10,12\}$ & $+1$ (QR) & $\bar\nu$ & $+1$ \\
$O_2=\{2,5,6\}$ & $-1$ (NQR) & Quarks & $+\tfrac{1}{3}$ \\
$O_4=\{7,8,11\}$ & $-1$ (NQR) & $e,\mu,\tau$ & $-1$ \\
\bottomrule
\end{tabular}
\end{center}
Beide Werte des Legendre-Symbols treten mit verschiedenen $Y$ auf; $Y$ ist keine
Funktion des Legendre-Symbols allein \SM.

\section{Zusammenfassung und Stand von Satz~E (Dok.~346)}

\begin{table}[H]
\centering
\caption{Ergebnisse Dok.~347}
\begin{tabular}{lll}
\toprule
Satz & Inhalt & Status \\
\midrule
A & $Y$ aus QR/NQR-Bit allein nicht festgelegt & \SM \\
B & $I_3 = \pm\tfrac{1}{2}$ aus jE7-Projektor (Su/Sd) & \BM \\
C & $Q = I_3 + Y/2$ mit $I_3$ aus Galois, $Y$ aus dem SM & \SM \\
D & Legendre-Symbol trennt $Y$ auf Orbit-Ebene nicht & \SM \\
\bottomrule
\end{tabular}
\end{table}

\textbf{Satz~E aus Dok.~346 [\SM] ist nur teilweise geschlossen.} Der Isospin $I_3$
folgt aus dem jE7-Projektor (Su/Sd-Trennung, Dok.~344) \BM; die Ladungsquantisierung
aus dem Legendre-Symbol mod~13 bleibt \BM\ (Dok.~346). Die Hyperladung $Y$ folgt
nicht aus dem QR/NQR-Bit allein und bleibt offen \SM.

\textbf{Was noch offen bleibt \SM:}
Die Hyperladung $Y$ aus $N\bmod3$ und der Teilchen/Antiteilchen-Unterscheidung (Satz~A), die explizite SU$(2)_L$-Dublett-Struktur (linkshändige Paare aus der Galois-Algebra)
und die CKM-Mischungswinkel aus Gal$(GF(27)/GF(3))=\mathbb{Z}_3$.


\section{Was noch fehlt: SU$(2)_L$ und CKM}

\subsection{SU$(2)_L$-Chiralität --- strukturell offen \SM}

SU$(2)_L$ wirkt nur auf \emph{linkshändige} Felder. Chiralität ist die
$\gamma_5$-Eigenschaft eines Spinors: linkshändig = Eigenwert $-1$,
rechtshändig = Eigenwert $+1$. Im Standardmodell bilden linkshändige Fermionen
Dubletts $(\nu_e, e^-)_L$ und $(u, d)_L$; rechtshändige Felder sind Singuletts.

\textbf{Was die Galois-Algebra liefert.}
Die Z$_2$-Parität in GF$(27)^* = \mathbb{Z}_2 \times \mathbb{Z}_{13}$
trennt die primitiven Klassen $f_i$ ($k < 13$) von den Negationen $g_i$ ($k \geq 13$).
Das ist eine Teilchen/Antiteilchen-Trennung, keine Chiralitätstrennung.
Alle vier primitiven Orbits $O_1,\ldots,O_4$ liegen auf derselben Seite
($k < 13$) und sind damit alle gleich-chiral.

\textbf{Was fehlt.}
$\gamma_5$ ist ein Element der Clifford-Algebra Cl$(6)$:
$\gamma_5 = e_1 e_2 e_3 e_4 e_5 e_6$ (Pseudoskalar).
In GF$(27)^*$ gibt es kein direktes Äquivalent.
Um die SU$(2)_L$-Dublett-Struktur explizit herzuleiten, bräuchte man
die Verbindung von GF$(27)^*$ zur Clifford-Algebra Cl$(6)$ ---
wo $\gamma_5$ und damit Chiralität definiert ist.
Das ist Fureysconstruction-Territorium (Dok.~344), nicht GF$(27)^*$-Territorium.

\textbf{Was nötig wäre.} Eine algebraische Brücke:
\[
  \text{GF}(27)^* \;\xrightarrow{?}\; \text{Cl}(6) \;\xrightarrow{\gamma_5}\; \text{Chiralität}
\]
Die erste Verbindung (GF$(27)^*$ in Cl$(6)$) ist in Dok.~341/344 begonnen,
aber noch nicht vollständig. Insbesondere ist die Einbettung der Galois-Orbits
in die Clifford-Gradstruktur (Grad~0--6) nicht explizit.

\subsection{CKM-Mischungswinkel --- konkreter Ansatz \SM}

Die CKM-Matrix beschreibt die Mischung zwischen up-artigen und down-artigen
Quarks. In FFGFT entsprechen die Generationen den drei Frobenius-Potenzen
$k \mapsto 3k$ innerhalb eines Orbits (Dok.~346, Satz~D). Die CKM-Mischung
ist dann die \emph{Überlappung zwischen verschiedenen Orbits über Generationen}.

\textbf{Was die Galois-Algebra liefert.}
GF$(27)$ trägt eine natürliche Spur-Bilinearform:
\[
  \langle a, b \rangle = \mathrm{Tr}_{\mathrm{GF}(27)/\mathrm{GF}(3)}(a \cdot b)
  = a \cdot b + (a \cdot b)^3 + (a \cdot b)^9.
\]
Diese innere Produktstruktur ermöglicht Überlappungen zwischen den Orbits.
Der Kabibbo-Winkel $\theta_C$ (zwischen 1.~und 2.~Generation) wäre:
\[
  \sin\theta_C \;=\; \frac{|\langle O_{2,\mathrm{Gen.1}},\, O_{4,\mathrm{Gen.2}}\rangle|}{
    \|O_{2,\mathrm{Gen.1}}\| \cdot \|O_{4,\mathrm{Gen.2}}\|}
\]
wobei $O_{2,\mathrm{Gen.1}} = g^2$ und $O_{4,\mathrm{Gen.2}} = g^{11}$ die
Orbit-Elemente der ersten bzw.\ zweiten Generation sind.

\textbf{Numerische Situation.}
Der experimentelle Wert ist $\sin\theta_C \approx 0{,}225$.
Ein erster Galois-Kandidat: $\sin(\pi/14) \approx 0{,}222$ (Abw.\ 1,2\,\%).
Eine algebraische Herleitung aus der Spur-Bilinearform ist jedoch noch nicht
vollständig ausgeführt.

\textbf{Was nötig wäre.}
\begin{enumerate}
\item Explizite Berechnung der Spur-Bilinearform $\langle g^k, g^j \rangle$
  für alle Orbit-Paare in GF$(27)$.
\item Identifikation der Generations-Zustände innerhalb der Orbits
  (welches Element = 1.~Gen., welches = 2.~Gen., welches = 3.~Gen.?).
\item Auswertung der Überlappungsmatrix als CKM-Kandidat.
\end{enumerate}
Das ist ein konkreter, ausführbarer algebraischer Weg. Er ist jedoch
nicht trivial und erfordert ein eigenes Dokument.

\textbf{Fazit.} SU$(2)_L$-Chiralität erfordert die Clifford-Brücke ---
kein GF$(27)^*$-Phänomen allein. CKM-Mischung erfordert die
Spur-Bilinearform-Struktur --- ein konkreter nächster Schritt, der
die Generationszuordnung innerhalb der Orbits voraussetzt.
Beide bleiben \SM, sind aber präzise formuliert und ausarbeitbar.

\section*{Prüfskript}
\texttt{pruef\_347\_gell\_mann.py} --- Orbit-Zuordnung nach Dok.~346, Nachweis, dass $Y$ keine Funktion des QR/NQR-Bits ist (Satz~A, D), $I_3$ (Satz~B) und $Q=I_3+Y/2$ mit SM-Werten (Satz~C); alle Assertions bestanden.

\section*{Vermerk zur Verwendung von KI-Assistenz}
Erstellt mit Unterstützung von Claude (Anthropic).
Physikalische Interpretationen und Schlussfolgerungen:
Johann Pascher (ORCID 0009-0000-6518-4064).

\begin{thebibliography}{9}
\bibitem{P344} J.~Pascher, \textit{Dok.~344: Furey-Fermionen in GALG und FFGFT}, FFGFT-Korpus (Sept.\ 2026).
\bibitem{P346} J.~Pascher, \textit{Dok.~346: Ladungsquantisierung aus GF$(27)^*$}, FFGFT-Korpus (Sept.\ 2026).
\bibitem{P339} J.~Pascher, \textit{Dok.~339: Frobenius-Trennung}, FFGFT-Korpus (Aug.\ 2026).
\bibitem{GellMann1956} M.~Gell-Mann, \textit{The interpretation of the new particles as displaced charge multiplets}, Nuovo Cimento~4 Suppl.~2 (1956) 848.
\bibitem{Nishijima1955} K.~Nishijima, \textit{Charge independence theory of V particles}, Prog.\ Theor.\ Phys.\ 13 (1955) 285.
\bibitem{PDG2022} R.L.~Workman \emph{et al.}\ (Particle Data Group),
Prog.\ Theor.\ Exp.\ Phys.\ \textbf{2022}, 083C01 (2022).
\bibitem{Furey2015} C.~Furey, Ph.D.\ Thesis, Waterloo 2015, arXiv:1611.09182.
\end{thebibliography}
